{\Large\textbf{Oscillating rope}}

A heavy uniform rope of length $L$ is suspended vertically from the ceiling.
The rope may oscillate around its equilibrium position with different natural
frequencies, which will be denoted by $f_i$ $(i = 1, 2, ...)$ in an ascending
order. The figure below presents the shape of the rope in the first three
natural vibrations, as obtained by a computer simulation. Note that the
horizontal and vertical scales in the figures are not equal. You may assume
that the actual lateral displacement of the rope is much smaller than its
length (small-amplitude approximation).

\begin{center}\includegraphics[width=0.56\linewidth]{figures/EuPhO_2017_Q1_fig1.png}\end{center}

Figura 1: Shapes of the oscillating rope for the first three natural vibrations
($i =$ 1,2,3... from left to right).

a) Develop a simplified model, which will allow you to estimate the frequency
$f_1$ of the first (fundamental) vibration of the rope. Hence, calculate
approximately $f_1$ for a rope of length $L = 1.0\ m$. (Assume that
$g = 9.8\ m/s$.)

b) Take necessary readings from the figure to estimate the frequency ratio
$f_1{:}f_2{:}f_3$.
